\documentclass[11pt]{article}
\usepackage{mathblatt}


\blattfuss{Wachstum}{Lösungen}
\begin{document}
\noindent{\large\bfseries Wachstum \textperiodcentered{} Lösungen}\par\medskip
\begin{pfloesung}{}
\lz{1.}{$55\,€$}{$50 \cdot 1{,}1$}{}
\end{pfloesung}
\begin{pfloesung}{}
\lz{2.}{$11$}{$y$}{}
\end{pfloesung}
\begin{pfloesung}{}
\lz{3.}{$6$}{}{}
\end{pfloesung}
\begin{pfloesung}{}
\lz{4.}{$30\,€$}{}{}
\end{pfloesung}
\begin{pfloesung}{}
\lz{5.}{$1{,}4$}{}{}
\end{pfloesung}
\begin{pfloesung}{}
\lz{6.}{$32$}{}{}
\end{pfloesung}
\begin{pfloesung}{}
\lz{7.}{der Quotient}{}{}
\end{pfloesung}
\begin{pfloesung}{}
\lz{8.}{die Differenz}{}{}
\end{pfloesung}
\begin{pfloesung}{}
\lz{9.}{$439{,}23$}{$300 \cdot 1{,}1^4$}{}
\end{pfloesung}
\begin{pfloesung}{}
\lz{10.}{der Startwert steht im Jahr $0$, er wird nicht noch einmal mit $1{,}1$ multipliziert}{$300$}{}
\end{pfloesung}
\begin{pfloesung}{}
\lz{11.}{$300$}{$375 : 1{,}25$}{}
\end{pfloesung}
\begin{pfloesung}{}
\lz{12.}{$287{,}5$}{$250 \cdot 1{,}15$}{}
\end{pfloesung}
\begin{pfloesung}{}
\lz{13.}{$343$}{$700$; $490$}{}
\end{pfloesung}
\begin{pfloesung}{}
\lz{14.}{$14\,739$}{$14\,739\,€$: $24\,000$; $20\,400$; $17\,340$}{}
\end{pfloesung}
\begin{pfloesung}{}
\lz{15.}{$1\,331$}{Jahr $3$, denn $1\,000 \cdot 1{,}1^3$}{}
\end{pfloesung}
\begin{pfloesung}{}
\lz{16.}{$484$}{$484$: $440 \cdot 1{,}1$; Probe: $484 \cdot 1{,}1 = 532{,}4$}{}
\end{pfloesung}
\begin{pfloesung}{{\small\color{mbgrau}P10 2020 · 4a}}
\lz{17.}{2019: 54; 2022: 57,3 · 1,03 = 59,019 $\approx$ 59,0}{57,3 · 1,03 = 59,019 $\approx$ 59,0}{}
\end{pfloesung}
\begin{pfloesung}{{\small\color{mbgrau}P10 2016 · 4a}}
\lz{18.}{4,45 mg; 4 h}{5,00 · 0,89 = 4,45}{}
\end{pfloesung}
\begin{pfloesung}{{\small\color{mbgrau}P10 2017 · 7a}}
\lz{19.}{756,9 $\approx$ 757 hPa; $\approx$ 328 hPa}{870 · 0,87 = 756,9 $\approx$ 757}{}
\end{pfloesung}
\begin{pfloesung}{{\small\color{mbgrau}P10 2019 · 7a}}
\lz{20.}{10 kg; 10,816 kg; $\approx$ 12,167 kg}{Wachstumsfaktor 1,04; 10,4 · 1,04 = 10,816}{}
\end{pfloesung}
\begin{pfloesung}{}
\lz{21.}{$350$}{}{}
\end{pfloesung}
\begin{pfloesung}{}
\lz{22.}{$225$}{}{}
\end{pfloesung}
\begin{pfloesung}{{\small\color{mbgrau}P10 2025 · 7a}}
\lz{23.}{Punkte (0|80), (1|112), (2|157), (3|220), (4|308) eingetragen, zunehmend steiler}{Skala wählen}{}
\end{pfloesung}
\begin{pfloesung}{{\small\color{mbgrau}P10 2021 · 6b}}
\lz{24.}{x-Achse 1 Kästchen = 5 min, y-Achse 1 Kästchen = 2 cm; fallende Strecke von (0; 40) bis (80; 24)}{80 min $\mathrel{\widehat{=}}$ 16 Kästchen; 40 cm $\mathrel{\widehat{=}}$ 20 Kästchen}{}
\end{pfloesung}
\begin{pfloesung}{}
\lz{25.}{$192$}{$50$; $98$; Ablesegenauigkeit $\pm 5$}{}
\end{pfloesung}
\begin{pfloesung}{}
\lz{26.}{Punkte $(0 | 100)$, $(1 | 120)$, $(2 | 144)$, $(3 | 173)$, $(4 | 207)$}{}{}
\end{pfloesung}
\begin{pfloesung}{}
\lz{27.}{$125$ liegt genau in der Mitte zwischen $100$ und $150$}{Punkte $(0 | 125)$, $(1 | 150)$, $(2 | 180)$, $(3 | 216)$}{}
\end{pfloesung}
\begin{pfloesung}{}
\lz{28.}{glatte, nach rechts steiler werdende Kurve durch $(0 | 100)$ bis $(4 | 207)$}{}{}
\end{pfloesung}
\begin{pfloesung}{}
\lz{29.}{Punkte eingetragen}{waagerecht $1$ Jahr je Kästchen, senkrecht $100\,€$ je Kästchen; $531$ passt in $800$}{}
\end{pfloesung}
\begin{pfloesung}{{\small\color{mbgrau}P10 2016 · 4b}}
\lz{30.}{x-Achse Zeit in h, 1 h je Teilstrich; y-Achse Masse in mg, 1 mg je Teilstrich; sieben Punkte, fallende Kurve}{Skala x 0 bis 8, y 0 bis 7}{}
\end{pfloesung}
\begin{pfloesung}{{\small\color{mbgrau}P10 2017 · 7b}}
\lz{31.}{x-Achse Höhe in km, 2 Kästchen je km; y-Achse Luftdruck in hPa, 1 Kästchen je 50 hPa; sieben Punkte, fallende Kurve}{x 0 bis 10 km, y 0 bis 1000 hPa}{}
\end{pfloesung}
\begin{pfloesung}{}
\lz{32.}{der Faktor ist größer als eins}{mehr}{}
\end{pfloesung}
\begin{pfloesung}{}
\lz{33.}{$2016$ ist der Schritt null}{$t = 5$}{}
\end{pfloesung}
\begin{pfloesung}{}
\lz{34.}{$1{,}4$}{}{}
\end{pfloesung}
\begin{pfloesung}{}
\lz{35.}{$1\,536$}{$2\,400$; $1\,920$}{}
\end{pfloesung}
\begin{pfloesung}{}
\lz{36.}{$40$ und $1{,}3$}{Anfangswert und Faktor}{}
\end{pfloesung}
\begin{pfloesung}{}
\lz{37.}{Abnahme um $10\,\%$ heißt Faktor $0{,}9$}{$B(t) = 2\,000 \cdot 0{,}9^t$}{}
\end{pfloesung}
\begin{pfloesung}{}
\lz{38.}{$600 \cdot 0{,}8^t$}{$m(t)$}{}
\end{pfloesung}
\begin{pfloesung}{}
\lz{39.}{$1\,536\,€$}{$3\,000 \cdot 0{,}8^3$}{}
\end{pfloesung}
\begin{pfloesung}{}
\lz{40.}{$0{,}9^2$, Faktor $0{,}9$}{$10\,\%$: $810 : 1\,000 = 0{,}81$}{}
\end{pfloesung}
\begin{pfloesung}{}
\lz{41.}{$40 \cdot 1{,}3^t$}{$A(t)$}{}
\end{pfloesung}
\begin{pfloesung}{}
\lz{42.}{$1{,}2$, also $20\,\%$ Zunahme}{$420 : 350$}{}
\end{pfloesung}
\begin{pfloesung}{}
\lz{43.}{Zunahme um $15\,\%$}{}{}
\end{pfloesung}
\begin{pfloesung}{}
\lz{44.}{Startwert $80$ und Quotient $1{,}25$}{$f(t) = 80 \cdot 1{,}25^t$}{}
\end{pfloesung}
\begin{pfloesung}{}
\lz{45.}{$0{,}75$}{}{}
\end{pfloesung}
\begin{pfloesung}{}
\lz{46.}{$1{,}05$: Faktor, jedes Jahr $5\,\%$ mehr}{$250$: Mitglieder zu Beginn}{}
\end{pfloesung}
\begin{pfloesung}{{\small\color{mbgrau}P10 2017 · 7c}}
\lz{47.}{y = 1000 · $0{,}87^{x}$}{}{}
\end{pfloesung}
\begin{pfloesung}{{\small\color{mbgrau}P10 2018 · 2a}}
\lz{48.}{600 000 € · 1,08 = 648 000 €}{48 000 € : 600 000 € = 8 \%}{}
\end{pfloesung}
\begin{pfloesung}{{\small\color{mbgrau}P10 2019 · 7b}}
\lz{49.}{10: Anfangsmasse in kg; 1,04: Wachstumsfaktor, +4 \% je Woche; x: Zeit in Wochen}{}{}
\end{pfloesung}
\begin{pfloesung}{}
\lz{50.}{der Bestand wächst jedes Jahr mit dem Faktor $1{,}2$, also um $20\,\%$}{$1{,}2$; $1{,}2$; $1{,}2$}{}
\end{pfloesung}
\begin{pfloesung}{}
\lz{51.}{$\approx 3\,546$}{$t = 6$: $2\,500 \cdot 1{,}06^6$}{}
\end{pfloesung}
\begin{pfloesung}{}
\lz{52.}{die Behauptung stimmt nicht, es fehlen knapp $40\,€$}{$2\,000 \cdot 1{,}04^{10} \approx 2\,960{,}49\,€$}{}
\end{pfloesung}
\begin{pfloesung}{{\small\color{mbgrau}P10 2017 · 7d}}
\lz{53.}{ja, 1000 · 0,87$^{4}$ $\approx$ 573 hPa}{0,87$^{4}$ $\approx$ 0,573}{}
\end{pfloesung}
\begin{pfloesung}{{\small\color{mbgrau}P10 2016 · 4e}}
\lz{54.}{$\approx$ 0,31 mg}{}{}
\end{pfloesung}
\begin{pfloesung}{{\small\color{mbgrau}P10 2020 · 4c}}
\lz{55.}{nein, Studie 1 54 · 1,03$^{14}$ $\approx$ 81,7 gegen Studie 2 = 75,6 (Tausend Liter)}{54 · 1,03$^{14}$ $\approx$ 81,7; 54 · 1,4 = 75,6}{}
\end{pfloesung}
\begin{pfloesung}{{\small\color{mbgrau}P10 2025 · 7b}}
\lz{56.}{112 : 80 = 1,4 (alle Quotienten $\approx$ 1,4); N(t) = 80 · $1{,}4^{t}$; wahr, N(24) $\approx$ 257 136 > 250 000}{$1{,}4^{24}$ $\approx$ 3214,2}{}
\end{pfloesung}
\begin{pfloesung}{}
\lz{57.}{$100$}{}{}
\end{pfloesung}
\begin{pfloesung}{}
\lz{58.}{nur er geht durch $(0 | 0)$}{$h$}{}
\end{pfloesung}
\begin{pfloesung}{}
\lz{59.}{fallende Kurve, die immer flacher wird}{$f$}{}
\end{pfloesung}
\begin{pfloesung}{}
\lz{60.}{$f$ passt: Start bei $100$, gekrümmt. $g$ passt nicht: Gerade. $h$ passt nicht: Start im Ursprung, die Fläche ist aber schon zu Beginn $100$}{}{}
\end{pfloesung}
\begin{pfloesung}{}
\lz{61.}{Start bei $300$ und Gerade, weil jede Minute derselbe Betrag dazukommt}{$g$}{}
\end{pfloesung}
\begin{pfloesung}{}
\lz{62.}{$f$ passt: Start bei $45$, Kurve. $g$ passt nicht: Gerade, also jede Woche gleich viel. $h$ passt nicht: Start bei $0$, das Kalb wiegt aber $45\,\mathrm{kg}$}{}{}
\end{pfloesung}
\begin{pfloesung}{}
\lz{63.}{$b$ ist die Gerade}{}{}
\end{pfloesung}
\begin{pfloesung}{}
\lz{64.}{eine Gerade hat immer denselben Zuwachs}{Gleicher Prozentsatz heißt gleicher Faktor: die Zuwächse werden von Jahr zu Jahr größer und der Graph steiler}{}
\end{pfloesung}
\begin{pfloesung}{{\small\color{mbgrau}P10 2019 · 7c}}
\lz{65.}{Graph 1 nein (linear); Graph 2 ja (Start 10 kg, Zuwächse wachsen); Graph 3 nein (Start bei 0 kg)}{Startwert 10 kg; gleicher Prozentsatz $\Rightarrow$ wachsende Zuwächse}{}
\end{pfloesung}
\begin{pfloesung}{\textbf{Prüfstein} \textperiodcentered{} P10 2026 FOR · 7}
\lz{a)}{2026: 650,00 €; 2029: 687,76 €}{Wachstumsfaktor 1,019; 674,93 · 1,019}{}
\lz{b)}{B; A: Miete beginnt bei 0 € statt bei 650 €; C: linear, jedes Jahr gleicher Betrag statt gleicher Prozentsatz}{Startwert 650 €; konstanter Prozentsatz $\Rightarrow$ wachsende Zuwächse}{}
\lz{c)}{M(t) = 650 · $1{,}019^{t}$; M(14) $\approx$ 845,96 € (rund 846 €)}{t = 2040 $-$ 2026 = 14; $1{,}019^{14}$ $\approx$ 1,3015}{}
\end{pfloesung}
\end{document}